\chapter{2023年全国甲卷(文)}

\section{单选题}

\begin{problem}
设全集 \(U=\{1,2,3,4,5\}\),集合 \(M=\{1,4\}\),\(N=\{2,5\}\),则 \(N\cup\complement_U M=\)(\quad)

\choices
    {\(\{2,3,5\}\)}
    {\(\{1,3,4\}\)}
    {\(\{1,2,4,5\}\)}
    {\(\{2,3,4,5\}\)}
\end{problem}

\begin{problem}
\(\dfrac{5\left(1+\i^3\right)}{\left(2+\i\right)\left(2-\i\right)}=\)(\quad)

\choices
    {\(-1\)}
    {\(1\)}
    {\(1-\i\)}
    {\(1+\i\)}
\end{problem}

\begin{problem}
已知向量 \(\bs{a}=\left(3,1\right)\),\(\bs{b}=\left(2,2\right)\),则 \(\cos\left(\bs{a}+\bs{b},\bs{a}-\bs{b}\right)=\)(\quad)

\choices
    {\(\dfrac{1}{17}\)}
    {\(\dfrac{\sqrt{17}}{17}\)}
    {\(\dfrac{\sqrt{5}}{5}\)}
    {\(\dfrac{2\sqrt{5}}{5}\)}
\end{problem}

\begin{problem}
某校文艺部有 4 名学生,其中高一、高二年级各 2 名. 从这 4 名学生中随机选 2 名组织校文艺汇演,则这 2 名学生来自不同年级的概率为(\quad)

\choices
    {\(\dfrac{1}{6}\)}
    {\(\dfrac{1}{3}\)}
    {\(\dfrac{1}{2}\)}
    {\(\dfrac{2}{3}\)}
\end{problem}

\begin{problem}
记 \(S_n\) 为等差数列 \(\{a_n\}\) 的前 \(n\) 项和. 若 \(a_2+a_6=10\),\(a_4a_8=45\),则 \(S_5=\)(\quad)

\choices
    {\(25\)}
    {\(22\)}
    {\(20\)}
    {\(15\)}
\end{problem}

\begin{problem}
执行下边的程序框图,则输出的 \(B=\)(\quad)

\bitmapfigure[max width=0.56\linewidth]{img/2023/national_paper_a_liberal/q6_fig1.png}

\choices
    {\(21\)}
    {\(34\)}
    {\(55\)}
    {\(89\)}
\end{problem}

\begin{problem}
设 \(F_1\),\(F_2\) 为椭圆 \(C:\dfrac{x^2}{5}+y^2=1\) 的两个焦点,点 \(P\) 在 \(C\) 上,若 \(\overrightarrow{PF_1}\cdot\overrightarrow{PF_2}=0\),则 \(\left|PF_1\right|\cdot\left|PF_2\right|=\)(\quad)

\choices
    {\(1\)}
    {\(2\)}
    {\(4\)}
    {\(5\)}
\end{problem}

\begin{problem}
曲线 \(y=\dfrac{\e^x}{x+1}\) 在点 \(\left(1,\dfrac{\e}{2}\right)\) 处的切线方程为(\quad)

\choices
    {\(y=\dfrac{\e}{4}x\)}
    {\(y=\dfrac{\e}{2}x\)}
    {\(y=\dfrac{\e}{4}x+\dfrac{\e}{4}\)}
    {\(y=\dfrac{\e}{2}x+\dfrac{3\e}{4}\)}
\end{problem}

\begin{problem}
已知双曲线 \(\dfrac{x^2}{a^2}-\dfrac{y^2}{b^2}=1\left(a>0,b>0\right)\) 的离心率为 \(\sqrt{5}\),其中一条渐近线与圆 \(\left(x-2\right)^2+\left(y-3\right)^2=1\) 交于 \(A\),\(B\) 两点,则 \(|AB|=\)(\quad)

\choices
    {\(\dfrac{\sqrt{5}}{5}\)}
    {\(\dfrac{2\sqrt{5}}{5}\)}
    {\(\dfrac{3\sqrt{5}}{5}\)}
    {\(\dfrac{4\sqrt{5}}{5}\)}
\end{problem}

\begin{problem}
在三棱锥 \(P-ABC\) 中,\(\triangle ABC\) 是边长为 \(2\) 的等边三角形,\(PA=PB=2\),\(PC=\sqrt{6}\),则该棱锥的体积为(\quad)

\choices
    {\(1\)}
    {\(\sqrt{3}\)}
    {\(2\)}
    {\(3\)}
\end{problem}

\begin{problem}
已知函数 \(f(x)=\e^{-\left(x-1\right)^2}\). 记 \(a=f\left(\dfrac{\sqrt{2}}{2}\right)\),\(b=f\left(\dfrac{\sqrt{3}}{2}\right)\),\(c=f\left(\dfrac{\sqrt{6}}{2}\right)\),则(\quad)

\choices
    {\(b>c>a\)}
    {\(b>a>c\)}
    {\(c>b>a\)}
    {\(c>a>b\)}
\end{problem}

\begin{problem}
函数 \(y=f(x)\) 的图象由 \(y=\cos\left(2x+\dfrac{\pi}{6}\right)\) 的图象向左平移 \(\dfrac{\pi}{6}\) 个单位长度得到,则 \(y=f(x)\) 的图象与直线 \(y=\dfrac{1}{2}x-\dfrac{1}{2}\) 的交点个数为(\quad)

\choices
    {\(1\)}
    {\(2\)}
    {\(3\)}
    {\(4\)}
\end{problem}

\section{填空题}

\begin{problem}
记 \(S_n\) 为等比数列 \(\{a_n\}\) 的前 \(n\) 项和. 若 \(8S_6=7S_3\),则 \(\{a_n\}\) 的公比为\fillinblank{}.
\end{problem}

\begin{problem}
若 \(f(x)=\left(x-1\right)^2+ax+\sin\left(x+\dfrac{\pi}{2}\right)\) 为偶函数,则 \(a=\fillinblank{}\).
\end{problem}

\begin{problem}
若 \(x\),\(y\) 满足约束条件
\(\begin{cases}
3x-2y\le3,\\
-2x+3y\le3,\\
x+y\ge1,
\end{cases}\)
则 \(z=3x+2y\) 的最大值为\fillinblank{}.
\end{problem}

\begin{problem}
在正方体 \(ABCD-A_1B_1C_1D_1\) 中,\(AB=4\),\(O\) 为 \(AC_1\) 的中点,若该正方体的棱与球 \(O\) 的球面有公共点,则球 \(O\) 的半径的取值范围是\fillinblank{}.
\end{problem}

\section{解答题}

\begin{problem}
记 \(\triangle ABC\) 的内角 \(A\),\(B\),\(C\) 的对边分别为 \(a\),\(b\),\(c\),已知 \(\dfrac{b^2+c^2-a^2}{\cos A}=2\).
\begin{enumerate}
\item 求 \(bc\);
\item 若 \(\dfrac{a\cos B-b\cos A}{a\cos B+b\cos A}-\dfrac{b}{c}=1\),求 \(\triangle ABC\) 面积.
\end{enumerate}
\end{problem}

\begin{problem}
如图,在三棱柱 \(ABC-A_1B_1C_1\) 中,\(A_1C\perp\) 平面 \(ABC\),\(\angle ACB=90^\circ\).

\begin{enumerate}
\item 证明:平面 \(ACC_1A_1\perp\) 平面 \(BCC_1B_1\);
\item 设 \(AB=A_1B\),\(AA_1=2\),求四棱锥 \(A_1-BB_1C_1C\) 的高.
\end{enumerate}

\bitmapfigure[max width=6.0cm]{img/2023/national_paper_a_liberal/q18_fig1.png}
\end{problem}

\begin{problem}
一项试验旨在研究臭氧效应,试验方案如下:选 40 只小白鼠,随机地将其中 20 只分配到试验组,另外 20 只分配到对照组,试验组的小白鼠饲养在高浓度臭氧环境,对照组的小白鼠饲养在正常环境,一段时间后统计每只小白鼠体重的增加量(单位:\(\mathrm{g}\)). 试验结果如下:

对照组的小白鼠体重的增加量从小到大排序为

15.2,18.8,20.2,21.3,22.5,23.2,25.8,26.5,27.5,30.1,

32.6,34.3,34.8,35.6,35.6,35.8,36.2,37.3,40.5,43.2

试验组的小白鼠体重的增加量从小到大排序为

7.8,9.2,11.4,12.4,13.2,15.5,16.5,18.0,18.8,19.2,

19.8,20.2,21.6,22.8,23.6,23.9,25.1,28.2,32.3,36.5
\begin{enumerate}
\item 计算试验组的样本平均数;
\item
\begin{enumerate}
\item 求 40 只小白鼠体重的增加量的中位数 \(m\),再分别统计两样本中小于 \(m\) 与不小于 \(m\) 的数据的个数,完成如下列联表
\begin{center}
\begin{tabular}{|c|c|c|}
\hline
& \(<m\) & \(\ge m\) \\
\hline
对照组 &  &  \\
\hline
试验组 &  &  \\
\hline
\end{tabular}
\end{center}
\item 根据列联表,能否有 \(95\%\) 的把握认为小白鼠在高浓度臭氧环境中与在正常环境中体重的增加量有差异?
\end{enumerate}
\end{enumerate}

附:\(K^2=\dfrac{n\left(ad-bc\right)^2}{\left(a+b\right)\left(c+d\right)\left(a+c\right)\left(b+d\right)}\),

\begin{center}
\begin{tabular}{c|ccc}
\hline
\(P\left(K^2\ge k\right)\) & 0.100 & 0.050 & 0.010 \\
\hline
\(k\) & 2.706 & 3.841 & 6.635 \\
\hline
\end{tabular}
\end{center}
\end{problem}

\begin{problem}
已知函数 \(f(x)=ax-\dfrac{\sin x}{\cos^2x}\),\(x\in\left(0,\dfrac{\pi}{2}\right)\).
\begin{enumerate}
\item 当 \(a=1\) 时,讨论 \(f(x)\) 的单调性;
\item 若 \(f(x)+\sin x<0\),求 \(a\) 的取值范围.
\end{enumerate}
\end{problem}

\begin{problem}
已知直线 \(x-2y+1=0\) 与抛物线 \(C:y^2=2px\left(p>0\right)\) 交于 \(A\),\(B\) 两点,\(|AB|=4\sqrt{15}\).
\begin{enumerate}
\item 求 \(p\);
\item 设 \(F\) 为 \(C\) 的焦点,\(M\),\(N\) 为 \(C\) 上两点,且 \(\overrightarrow{FM}\cdot\overrightarrow{FN}=0\),求 \(\triangle MFN\) 面积的最小值.
\end{enumerate}
\end{problem}

\section{选考题:解答题}

\begin{problem}
已知点 \(P(2,1)\),直线 \(l\) 的参数方程为
\(\begin{cases}
x=2+t\cos\alpha,\\
y=1+t\sin\alpha,
\end{cases}\)(\(t\) 为参数),\(\alpha\) 为 \(l\) 的倾斜角,\(l\) 与 \(x\) 轴正半轴,\(y\) 轴正半轴分别交于 \(A\),\(B\),且 \(|PA|\cdot|PB|=4\).
\begin{enumerate}
\item 求 \(\alpha\);
\item 以坐标原点为极点,\(x\) 轴正半轴为极轴建立极坐标系,求 \(l\) 的极坐标方程.
\end{enumerate}
\end{problem}

\begin{problem}
已知 \(f(x)=2|x-a|-a\),\(a>0\).
\begin{enumerate}
\item 求不等式 \(f(x)<x\) 的解集;
\item 若曲线 \(y=f(x)\) 与 \(x\) 轴所围成的图形的面积为 \(2\),求 \(a\).
\end{enumerate}
\end{problem}

